\scanpage{131}
\mspage{355}

l'ens. des \unsure{pts} fixes de $U_{0,3}$ sous $\dot{\sigma}_\infty$ est réduit à un seul élément, ce qui implique aussi que
\deleted{\illegible{1}}
\[
\Pi_{0,3}^{\sigma_\infty}=\{1\}.
\]

Nous admettrons \unsure{encore} que ces faits \unsure{s'étendent} aux éléments de $\widehat{\Pi}_{0,3}$ --- ce qui sera le cas si (en plus de ce qui en a été admis à propos de $\rho$) on admet qu'on a \unsure{encore}
\[
\tag{38}
\widehat{\Pi}_{0,3}.
\]

Notons qu'on a ici
\[
\tag{39}
G^{\sigma_\infty}=\mathcal T_{0,3}^{\sigma_\infty}\simeq\mathbb Z/2\mathbb Z\times\mathbb Z/2\mathbb Z,
\]
sous-groupe engendré par $\dot{\sigma}_\infty,\unsure{\dot\tau|\tau}$, dont les éléments non \unsure{nuls} sont $\dot{\sigma}_\infty,\tau,\dot{\tau}'_\infty$ --- qui se \unsure{remontent} en $\sigma_\infty,\tau_\infty,\tau'_\infty$ ; et on a
\[
\tag{40}
\mathcal E^{\sigma_\infty}=\mathcal G_{0,3}^{\sigma_\infty}\simeq\mathbb Z/2\mathbb Z\times\mathbb Z/2\mathbb Z,
\]
groupe engendré par $\sigma_\infty,\tau_\infty$ (l'autre élément non trivial étant $\tau'_\infty$). Si donc on écrit une solution $h\inserted{\in\pi}$ (de (35) sous la forme (36)), avec $\beta\in\mathcal G_{0,3}$ seulement (pas seulement $\beta\in\pi$), alors l'indétermination est dans le groupe (40) d'ordre (4) ; De façon plus précise, si $\beta_0$ est la solution dans $\pi$, alors les autres solutions sont de la forme
\[
\tag{41}
\beta=\beta_0\sigma_\infty^\varepsilon\tau_\infty^{\varepsilon'}
\qquad \varepsilon,\varepsilon'\in\mathbb Z/2\mathbb Z
\]

\margin{130}

\scanpage{132}

(avec $\varepsilon,\varepsilon'$ uniquement déterminés). Pour qu'on ait $\beta\in\mathcal G_{0,3}^{+}$, il f. et s. que $\varepsilon'=0$ !

\margin{131}

\scanpage{133}
\mspage{356}

\margin{Ces notations seront \unsure{revisées} encore une fois au \S\ 46}

\uline{44}\qquad \uline{Retour sur les notations.} \uline{Mise en équations} \uline{dans} $\mathrm{Sl}(2,\widehat{\mathbb Z})$.

Soit $u$ un autom. de $\widehat{\Pi}_{0,3}$ \inserted{commutant extérieurement : $\unsure{\mathcal G_3}$\deleted{\illegible{1}}} --- ou ce qui revient au $\overline{\mathrm m}$, un autom. de $\widehat{\mathcal G}_{0,3}^{+}$ qui induit l'identité dans $\mathcal T_{0,3}^{+}=\unsure{\mathcal G_3}$. Alors $u$ est connu quand on connaît ses actions sur les deux générateurs $\sigma_\infty,\rho$ de $\widehat{\mathcal G}_{0,3}^{+}$, qui seront resp. des élts de $\widehat{\mathcal G}_{0,3}^{+}$ au dessus de $\dot{\sigma}_\infty$ et de $\dot\rho$,

soit
\[
\tag{1}
\boxed{u(\sigma_\infty)=h\sigma_\infty,\qquad u(\rho)=g\rho}
\qquad (h,g\in\widehat{\Pi}_{0,3})
\]
où $h,g$ \uline{sont assujettis aux conditions}
\[
\tag{2}
\boxed{h\sigma_\infty(h)=1\qquad g\rho(g)\rho^2(g)=1},
\]
exprimant respectivement les relations
\[
\tag{3}
\boxed{(h\sigma_\infty)^2=(g\rho)^3=1}.
\]

(Il y a en plus à vérifier la condition \unsure{liminaire} --- savoir que (3) sont des relations génératrices de $\widehat{\mathcal G}_{0,3}^{+}$ \dots).

\deleted{\illegible{1}}\inserted{Pour le moment}, il n'y a pas lieu d'imposer aucune relation entre $g$ et $h$ --- sauf les relations (2) sur chacun de ces éléments --- puisque $\widehat{\mathcal G}_{\unsure{0,3}}^{+}$ est le groupe ayant les relations génératrices
\[
\sigma_\infty^2=\rho^3=1.
\]

\margin{132}

\scanpage{134}

Les éléments $g$ et $h$, donc pris \unsure{comme} des autom. de $\widehat{\Pi}_{0,3}$, expriment les relations de commutation de\deleted{s} $u$ avec \deleted{$\sigma_\infty$} $\rho$ et $\sigma_\infty$ respectivement :
\[
\tag{4}
[u,\rho]=\underbrace{u\rho u^{-1}}_{u(\rho)}\rho^{-1}=\operatorname{int}(g),
\qquad
[u,\sigma_\infty]=u\sigma_\infty u^{-1}\sigma_\infty=\operatorname{int}(h)
\]
--- c'est une chose remarquable que la connaissance des éléments $g,h$ de $\pi$ qui vérifient (4), \uline{\unsure{implique|impliquent}} déjà la connaissance de $u$ (\unsure{ceci|ce qui} \unsure{au} fait que $\sigma_\infty,\rho$ engendrent $\mathcal G_{0,3}^{+}$ \dots).

Écrivons maintenant la compatibilité de $u$ avec la structure à lacets, \deleted{\illegible{3}}\inserted{On l'avait} \deleted{\unsure{formulé}}\inserted{\unsure{exprimé|exprimée}} précédemment par une relation de la forme
\[
u(\varepsilon_0)=\operatorname{int}(g_0)(\varepsilon_0^p)
\quad\text{ou}\quad
u(\ell_0)=\operatorname{int}(g_0)\ell_0^p,
\]
pour
\[
p\in\widehat{\mathbb Z}^{*},\qquad g_0\in\widehat{\Pi}_{0,3}
\]
convenables. La notation $g_0$ n'est pas tellement heureuse, puisque $g$ et $g_0$ jouent des rôles tout à fait dissemblables. L'indice $0$ dans $g_0$ était destiné à pouvoir écrire, plus généralement
\[
\deleted{\unsure{g_i}}\quad
u(\ell_i)=\operatorname{int}(g_i)\ell_i^p
\qquad i\in\{0,1,\infty\}.
\]

\margin{133}

\scanpage{135}
\mspage{357}

Je préfère utiliser la lettre $f_0$ --- quitte à supprimer l'indice s'il n'y a pas de confusion possible. La lettre $p$ pour un multiplicateur n'est pas \inserted{un choix} très heureux, car par la suite \deleted{\unsure{$p$}} on aura \inserted{\unsure{sûrement}} besoin de cette lettre pour désigner une caractéristique --- je préfère la lettre $\mu$ (initiale grecque de ``multiplicateur''), et la relation de lacets prend\inserted{\unsure{rait}} donc la forme
\[
\deleted{\unsure{(5)}}\qquad
u(\varepsilon_0)=\operatorname{int}(f_0)\varepsilon_0^\mu.
\qquad
\deleted{\unsure{\text{ou}}\quad u(\ell_0)=\operatorname{int}(f_0)}
\]

En fait, je préfère y remplacer $f_0$ par $f_0^{-1}$, pour obtenir par la suite des formules où la symétrie des rôles joués par $\rho$ et par $\sigma_\infty$ \unsure{apparaisse} plus nettement. --- notamment je voudrais que la formule de définition de $h_0$ en termes de $g_0$ (\S\ 43, formule (65,)
\[
h_0=g_0^{-1}\sigma_\infty(g_0)
\]
devienne plutôt
\[
h_0=g_0\sigma_\infty(g_0)^{-1},
\]
en harmonie avec
\[
\xi=\lambda\rho(\lambda)^{-1}\quad\dots
\]

Donc la ``relation de \unsure{lacets|lacet}'' prend la forme
\[
\tag{5}
u(\varepsilon_0)=\operatorname{int}(f_0^{-1})\varepsilon_0^\mu
\qquad f_0\in\widehat{\Pi}_{0,3},\ \mu\in\widehat{\mathbb Z}^{*},
\]
\deleted{\unsure{plutôt} \unsure{on} déduit aussitôt}
\[
\tag{6}
u(\ell_0)=\operatorname{int}(f_0^{-1})\ell_0^\mu,
\]
dont on déduit aussitôt

\margin{134}

\scanpage{136}

\[
\tag{7}
\left\{
\begin{aligned}
u(\varepsilon_1)&=\operatorname{int}(f_1^{-1})(\varepsilon_1^\mu),
&u(\ell_1)&=\operatorname{int}(f_1^{-1})(\ell_1^\mu),\\
u(\varepsilon_\infty)&=\operatorname{int}(f_\infty^{-1})(\varepsilon_\infty^\mu),
&u(\ell_\infty)&=\operatorname{int}(f_\infty^{-1})(\ell_\infty^\mu),
\end{aligned}
\right.
\]
\unsure{où|ou}
\[
\tag{8}
f_1=\rho(f_0)g^{-1},\qquad
f_\infty=\rho(f_1)g^{-1}=\rho^2(f_0)\rho(g)^{-1}g^{-1}.
\]

Dans la suite, on écrira simplement $f$ au lieu de $f_0$, et on va transformer l'équation (5) \deleted{\illegible{1}} en
\[
\tag{8bis}
\operatorname{int}(f)u(\varepsilon_0)=\varepsilon_0^\mu,
\]
en explicitant le premier membre
\[
\begin{aligned}
\operatorname{int}(f)u(\varepsilon_0)
&=f\,u(\sigma_\infty\rho)f^{-1}
=f\underbrace{u(\sigma_\infty)}_{h\sigma_\infty}\underbrace{u(\rho)}_{g\rho}f^{-1}
=fh\sigma_\infty g\rho f^{-1}\\
&=fh\sigma_\infty g\sigma_\infty^{-1}\underbrace{\sigma_\infty\rho}_{\varepsilon_0}f^{-1}
=fh\sigma_\infty(g)\underbrace{\varepsilon_0f^{-1}\varepsilon_0^{-1}}\varepsilon_0
\end{aligned}
\]
i.e.
\[
\tag{9}
\operatorname{int}(f)u(\varepsilon_0)
=\bigl[fh\sigma_\infty(g)(\sigma_\infty\rho)(f)^{-1}\bigr]\varepsilon_0
=\bigl[f,h\sigma_\infty(g\rho(f)^{-1}\bigr]\varepsilon_0
\]
où les \unsure{quatre|trois} facteurs $f,h,\sigma_\infty(g\rho(f)^{-1}$ dans $[\ ]$ sont $\in\widehat{\Pi}_{0,3}$. Ainsi \unsure{(5)|(9)} ou (8), prend la forme
\[
fh\sigma_\infty(g\rho(f)^{-1})=\varepsilon_0^{\mu-1}\,\deleted{\illegible{1}}
\]
ce qui donne, en posant
\[
\tag{10}
\mu-1=2\nu\qquad(\nu\in\widehat{\mathbb Z})
\]
(on note donc $\nu$ ce qui était désigné par $\gamma$ précédemment) on trouve la relation
\[
fh\sigma_\infty\bigl(g\rho(f)^{-1}\bigr)=\ell_0^\nu
\]
ou encore, en transformant par $\sigma_\infty$
\[
\sigma_\infty(\unsure{f})\sigma_\infty(h)g\rho(f)^{-1}=\ell_1^\nu
\]
; soit

\margin{135}

\scanpage{137}
\mspage{358}

\[
\boxed{\sigma_\infty(h).g=\sigma_\infty(f)^{-1}\ell_1^\nu\rho(f)}.
\qquad (11)
\]

Cette relation, qui relie $h$ et $g$ via $f$, peut à juste titre s'appeler la ``relation lacets''.
On voit, par cette relation, que \deleted{\illegible{1}} non seulement la connaissance \inserted{de $u$ i.e.} \unsure{des|de} $h$ et $g$ détermine $f$ et $\mu$ i.e. $f$ et $\nu$, mais aussi inversement, une fois $f$ et $\nu$ connus, $g$ et $h$ se déterminent mutuellement via (11).

\margin{$f$ déterminé mod $L_0$ seulement $f\mapsto\ell_0^n f$ $(n\in\widehat{\mathbb Z})$.}

Notons que pour $u$ correspondant à : \deleted{\unsure{$f$|$\xi$}}, $g,h,f$, si $\gamma\in\widehat{\Pi}_{0,3}$, l'\unsure{autom.} \inserted{$u'=$} $\operatorname{int}(\gamma)\circ u$ correspond à
\[
g'=\gamma g\rho(\gamma^{-1}),\quad
h'=\gamma h\sigma_\infty(\gamma)^{-1},\quad
f'=f\gamma^{-1}.
\qquad (12)
\]

\margin{Cela vaut plus \unsure{généralement} si $\gamma\in\unsure{\widetilde{\mathcal G}_{0,3}|\mathcal G_{0,3}}$, et même pour $\gamma\in\unsure{\widehat{\widetilde{\mathcal G}}_{0,3}|\widehat{\mathcal G}_{0,3}}$.}

Les équations (2) ``en $h$'' et ``en $g$'' se résolvent paramétriquement de façon unique par
\[
g=\alpha\rho(\alpha)^{-1},\quad
h=\beta\sigma_\infty(\beta)^{-1}
\qquad (13)
\]
où
\[
\left\{\begin{array}{l}
\alpha\in\widehat{\widetilde{\Pi}}_{0,3}=\widehat{\Pi}_{0,3}.\{1,\tau_\infty\}\\
\beta\in\widehat{\Pi}_{0,3}
\end{array}\right.
\qquad (14)
\]
de façon précise, la donnée des $g,h$ satisfaisant (2) équivaut resp. à celle des $\alpha,\beta$ dans $\widehat{\widetilde{\Pi}}_{0,3}$ et $\widehat{\Pi}_{0,3}$, via les formules (13).

\margin{136}

\scanpage{138}

\inserted{Les éléments} $g$ et $h$ étant ainsi explicités en termes des $\alpha$ et $\beta$, la relation \inserted{des lacets} (11) peut être considérée comme une relation entre $\alpha,\beta,f,\nu$
\[
\sigma_\infty(\beta\sigma_\infty(\beta)^{-1})(\alpha\rho(\alpha)^{-1})
=\sigma_\infty(f)^{-1}\ell_1^\nu\rho(f),
\qquad (15)
\]
où \deleted{$\nu,f$ \unsure{sont}} $\nu$ est déterminé de façon unique par $\alpha,\beta$, et $f$ est déterminé dans $L_0\backslash\widehat{\Pi}_{0,3}$.
(On notera que le deuxième membre de (15) ne change pas, par le changement de $f$ en $\ell_0^n f$, avec $n\in\widehat{\mathbb Z}$).

La dissymétrie principale, me semble-t-il, entre $\sigma_\infty$ et $\rho$ dans cette théorie, c'est que \deleted{\illegible{1}} l'on a tjrs $\beta\in\widehat{\Pi}_{0,3}$, tandis que $\alpha$ n'est pas nécessairement dans $\widehat{\Pi}_{0,3}$ -- il l'est, \unsure{on l'a vu}, ssi $\nu\equiv0\ (3)$ i.e. $\mu\equiv1\ (3)$.

\deleted{Donc si on utilise \illegible{1} $\operatorname{int}(\beta^{-1})$ \illegible{1}}
L'autom.
\[
u'=\operatorname{int}(\beta^{-1})u
\]
sera donc \unsure{congru} à $u$ mod autom \unsure{intérieurs}, et correspondra à
\[
\left\{\begin{array}{l}
g'=\beta^{-1}g\rho(\beta)=(\beta^{-1}\alpha)\rho(\beta^{-1}\alpha)^{-1}\\
h'=\beta^{-1}h\sigma_\infty(\beta)=1\\
f'=f\beta
\end{array}\right.
\qquad (16)
\]
-- ainsi $u$ est \deleted{conjugué} \inserted{congru mod} autom. int.

\margin{137}

\scanpage{139}
\mspage{359}

à un unique $u'$ tel que $u'(\sigma_\infty)=\sigma_\infty$ i.e. un $u'$ qui commute à $\sigma_\infty$. Mais dans \unsure{la notation symétrique|les notations symétriques} on prendra
\[
u''=\operatorname{int}(\alpha^{-1})u
\]
correspondant à
\[
\left\{\begin{array}{l}
g''=\alpha^{-1}g\rho(\alpha)=1\\
h''=\alpha^{-1}h\sigma_\infty(\alpha)=(\alpha^{-1}\beta)\sigma_\infty(\alpha^{-1}\beta)^{-1}\\
f''=f\alpha
\end{array}\right.
\qquad (17)
\]
mais $u''$ sera \inserted{seulement $\widehat{\widetilde{\Pi}}_{0,3}$-}conjugué \deleted{\unsure{par}} de $\unsure{u''|u}$ -- il sera \inserted{$\widehat{\Pi}_{0,3}$-}conjugué à un \deleted{élément de} \inserted{\unsure{autom}} qui commute à $\rho$ ssi $g$ est $\Pi_{0,3}$-$\rho$-conjugué à $1$ i.e ssi $\alpha\in\widehat{\Pi}_{0,3}$ i.e. ssi $\mu\equiv1\ (3)$ ou encore $\nu\equiv0\ (3)$.

\deleted{Quant au} Les calculs des \S\S\ précédents étaient relatifs au cas où $h=1$, i.e. $u$ commute à $\sigma_\infty$ -- donc \illegible{1} il reste les quantités $g$ et $f$, \deleted{on voit d'ailleurs} \inserted{plus le multiplicateur $\mu$}
\deleted{\inserted{le multiplicateur} pour $\mu$, \unsure{ou pour} $\nu=\frac{\mu-1}{2}$, est déterminé sans ambiguïté par}
-- une fois $\mu$ connu i.e $\nu$ connu, $g$ et $f$ se déterminent mutuellement par (11) ($f$ demeurant élément de $L_0\backslash\widehat{\Pi}_{0,3}$).

\margin{On verra que $\mu$ est déterminé déjà par $f$, \unsure{du moins modulo} signe, cf plus bas.}

Pour $f$ fixé, l'équation en $g$ peut s'interpréter comme une équation en $f$, ou mieux, en introduisant
\[
\varphi=f\sigma_\infty(f)^{-1}
\qquad (12)
\]

\margin{Cette notation n'est \unsure{utile|valable} que si $\beta=1$ \illegible{2}}
\margin{c'est essentiellement l'élément noté $h_0$ précédemment}
\margin{138}

\scanpage{140}

une ``équation en $\varphi$'' --
\[
(\ell_0^\nu\varphi)\rho(\ell_0^\nu\varphi)\rho^2(\ell_0^\nu\varphi)=1,
\qquad (13)
\]
ou encore, posant
\[
\boxed{\xi=\ell_0^\nu\varphi\quad\text{i.e. }\varphi=\ell_0^{-\nu}\xi}
\qquad (14)
\]
on peut l'interpréter comme ``équation en $\xi$''
\[
\xi\rho(\xi)\rho^2(\xi)=1.
\qquad (15)
\]
Ainsi, les autom. à lacets de $\widehat{\Pi}_{0,3}$ qui commutent à $\sigma_\infty$, \inserted{au multiplicateur $\mu=2\nu+1$,} correspondent aux couples de solutions $(\xi,\varphi)$ des équations
\[
\boxed{\varphi\sigma_\infty(\varphi)=1\qquad
\xi\rho(\xi)\rho^2(\xi)=1}
\qquad (16)
\]

\margin{On \unsure{trouvera|trouve} que \unsure{tant} $\xi$ que $\varphi$ \unsure{déterminent} déjà $\mu$, donc $\nu$.}

satisfaisant de plus la condition de compatibilité (14) (particulièrement simple \deleted{\unsure{(17)}} si $\mu=1$ i.e. \deleted{$\mu=$} $\nu=0$, puisqu'alors $\varphi=\xi$ !). On récupère $g$ (donc $u$) \deleted{par la} en fonction
\deleted{Inversement, si on \illegible{3}}
de $(\varphi,\xi)$ par la formule des lacets (11), qui ici prend la forme
\[
g=\sigma_\infty(f)^{-1}\ell_1^\nu\rho(f)
=(\ell_0^\nu f)^{-1}\bigl(\underbrace{\ell_0^\nu\varphi}_{\xi}\bigr)\rho(\ell_0^\nu f)
\qquad (17)
\]
$f\in\widehat{\Pi}_{0,3}$ étant déterminé de façon unique en termes de $\varphi$ par la formule (12).

De façon symétrique, on va s'intéresser aux autom à lacets $u$ de $\Pi_{0,3}$ qui commu-

\margin{139}

\scanpage{141}
\mspage{360}

tent à $\rho$, i.e. ceux pour lesquels on a $g=1$, décrits par suite en termes de $h,f$,
\[
u(\sigma_\infty)=h\sigma_\infty,\quad u(\rho)=1_{\deleted{\unsure{\rho}}}\qquad (18)
\]
\[
u(\varepsilon_0)=\operatorname{int}(f^{-1})(\varepsilon_0^\mu)
\qquad (19)
\]
cette dernière relation prenant la forme (11)
\[
\sigma_\infty(h)=\sigma_\infty(f)^{-1}\ell_1^\nu\rho(f),
\qquad (20)
\]
qui s'écrit aussi, en appliquant $\sigma_\infty$
\[
h=f^{-1}.\ell_0^\nu.(\sigma_\infty\rho)(f)\unsure{\underset{\cdot}{=}|=}
\qquad (21)
\]
``L'équation en $h$'' (2) \unsure{la}
\[
h\sigma_\infty(h)=1
\qquad \unsure{(*)}
\]
s'écrit donc comme une ``équation en $f$'', \deleted{qui} \inserted{avec}
\[
h\sigma_\infty(h)=f^{-1}\ell_0^\nu
\underbrace{(\sigma_\infty\rho)(f)\sigma_\infty(f)^{-1}}_{\sigma_\infty(\underbrace{\rho(f)f^{-1}}_{\psi^{-1}})}
\ell_1^\nu\rho(f)\deleted{\illegible{1}}
\]
donc \unsure{ça} s'écrit, en termes de $\psi=f\rho(f)^{-1}$
\[
\underbrace{\ell_0^\nu\sigma_\infty(\psi^{-1})}\ell_1^\nu
\underbrace{\rho(f)f^{-1}}_{\psi^{-1}}=1
\]
soit (en passant à l'inverse)
\[
\psi\ell_1^{-\nu}\underbrace{\sigma_\infty(\psi)\ell_0^{-\nu}}_{\sigma_\infty(\psi\ell_1^{-\nu})}=1
\quad\text{soit}\quad
\underbrace{\ell_0^{-\nu}\psi}_{\eta}
\underbrace{\ell_1^{-\nu}\sigma_\infty(\psi)}_{\sigma_\infty(\underbrace{\ell_0^{-\nu}\psi}_{\eta})}=1.
\]
\deleted{soit}
Donc, posant
\[
\psi\overset{\mathrm{dfn}}{=}f\rho(f)^{-1}
\qquad (22)
\]
\[
\boxed{\eta=\ell_0^{-\nu}\psi\quad\text{i.e. }\psi=\ell_0^\nu\eta}
\qquad (23)
\]
, on trouve que les automorphismes à lacets $u$ de $\widehat{\Pi}_{0,3}$ qui commutent à $\rho$ correspondent aux triples $(\eta,\psi,\nu)$, $\eta,\psi\in\widehat{\Pi}_{0,3}$, $\nu\in\widehat{\mathbb Z}$ (avec $2\nu+1\in\widehat{\mathbb Z}^{*}$), satisfaisant

\margin{140}

\scanpage{142}

\[
\boxed{\psi\rho(\psi)\rho^2(\psi)=1,\qquad\eta\sigma_\infty(\eta)=1},
\qquad (24)
\]
\margin{On \unsure{trouvera} \illegible{1} \unsure{encore}, \unsure{par} \unsure{tant} $\psi$ que $\eta$}

et en plus la relation de compatibilité (23) -- particulièrement simple si $\mu=1$ i.e. si $\nu=0$, car alors elle signifie $\psi=\eta$ ! Mais en plus de ces relations, il faut supposer que la solution $\psi$ de l'équation (24) ``en $\psi$'' \unsure{appartient|appartienne} à la famille principale, i.e. que, quand on écrit $\psi$ sous forme paramétrique
\[
\psi=f\rho(f)^{-1}
\qquad (22)
\]
avec $f\in\widehat{\widetilde{\Pi}}_{0,3}=\widehat{\Pi}_{0,3}.\{1,\tau_\infty\}$, on ait bien $f\in\widehat{\Pi}_{0,3}$ (donc que l'image de $f$ dans $\widehat{\widetilde{\Pi}}_{0,3}/\widehat{\Pi}_{0,3}\simeq\{1,\tau_\infty\}\simeq\mathbb Z/2$ soit nulle\dots).

Ayant ainsi déterminé $f$ en \uline{fonction} de $\psi$, \unsure{déterminer|détermine} $h$ par (21), d'où $u$ \deleted{\illegible{3}}.

\uline{Remarques} 1) Si cette condition supplémentaire sur $\psi$ n'est pas satisfaite, la \deleted{\unsure{détermination}} formule (21) garde un sens et définit $h\in\widehat{\Pi}_{0,3}$, satisfaisant ``l'équation en $h$'' $h\sigma_\infty(h)=1$ donc définissant un autom $u$ de $\widehat{\Pi}_{0,3}$, commutant à $\rho$.

On aura alors
\[
u(\ell_0)=\operatorname{int}(f^{-1})(\ell_0^\mu)
\]
ou encore, posant
\[
f=\inserted{\tau_\infty}f_0\deleted{\tau_\infty}\qquad\text{avec }f_0\in\widehat{\Pi}_{0,3}
\]
\[
u(\ell_0)=\operatorname{int}(f_0^{-1})(\tau_\infty(\ell_0^\mu))=\operatorname{int}(f_0^{-1})(\ell_0^{-\mu})
\]

\margin{141}

\scanpage{143}
\mspage{361}

i.e. $u$ est un automorphisme à lacets, de multiplicateur $-\mu=-(2\nu+1)$.

2) Notons que l'indétermination de $f$ (pour $u$ fixé, i.e. $g,h$ fixés) est dans
\[
 f\longmapsto f'=\ell_0^n f\;\deleted{\unsure{f}}
 \qquad(n\in\widehat{\mathbb Z})
 \tag{25}
\]
qui donne
\[
 \left\{\begin{aligned}
 \varphi&\longmapsto\varphi'=\ell_0^n\varphi\ell_1^{-n},&
 \psi&\longmapsto\psi'=\ell_0^n\psi\ell_1^{-n},\\
 \xi&\longmapsto\xi'=\ell_0^n\xi\ell_1^{-n},&
 \eta&\longmapsto\eta'=\ell_0^n\eta\ell_1^{-n},
 \end{aligned}\right.
 \tag{26}
\]
\margin{\uline{NB} On a de même
$\alpha_0\longmapsto\alpha_0'=\ell_0^n\alpha_0$
$\beta_0\longmapsto\beta_0'=\ell_0^n\beta_0$.}

(NB on a $\ell_1=\sigma_\infty(\ell_0)=\rho(\ell_0)$, donc $\xi'$ et $\psi'$ sont les $\rho$-transformés de $\xi,\psi$ par $\unsure{\ell_0^{-n}|\ell_0^n}$, et $\varphi',\eta'$ les $\sigma_\infty$-transformés de $\varphi,\eta$ par $\unsure{\ell_0^{-n}|\ell_0^n}$ --- en particulier les équations (16) (en $\varphi$ et en $\xi$) et (24) (en $\psi$ et en $\eta$) sont bien invariantes par ces transformations, ainsi d'ailleurs que les relations de compatibilité (14) resp. (23).

3) On a les formules suivantes \inserted{(réécriture de (11))}
\[
 \sigma_\infty(h)g
 = (\ell_0^\nu f)^{-1}\xi\rho(\ell_0^\nu f)
 = \sigma_\infty(f^{-1}\ell_0^\nu)\eta(f^{-1}\ell_0^\nu)^{-1}
 \tag{27}
\]
\deleted{(cas $h=1$)} --- exprimant $\inserted{\sigma_\infty(h)}g$ comme le $\rho$-conjugué de $\xi$ par $(\ell_0^\nu f)^{-1}$, et de façon symétrique \deleted{dans le cas $g=1$,} comme le $\sigma_\infty$-conjugué de $\eta$ par $\sigma_\infty(f^{-1}\ell_0^\nu)$. Si $h=1$, c'est la première expression qui est intéressante (impliquant l'équation en $\xi$), pour $g=1$, c'est la seconde, qu'on peut aussi réécrire \inserted{(en appliquant $\sigma_\infty$ aux 2 membres)}
\[
 h=(f^{-1}\ell_0^\nu)(\eta^{-1})\sigma_\infty(f^{-1}\ell_0^\nu)^{-1}
 \tag{28}
\]
\margin{142}

\scanpage{144}

exprimant $h$ comme $\sigma_\infty$-transformé de $\eta^{-1}$ par $f^{-1}\ell_0^\nu$. Si on veut l'exprimer comme \inserted{$\sigma_\infty$-}transformé de $\eta$ (au lieu de $\inserted{\sigma_\infty(\eta)=}\eta^{-1}$), on écrit
\[
 \deleted{\unsure{\eta^{-1}=\eta^{-1}\eta\sigma_\infty(\eta^{-1})}}
\]
\[
 \eta^{-1}=\sigma_\infty(\eta)=\eta^{-1}.\eta.\sigma_\infty(\eta)
\]
d'où
\[
 h=(f^{-1}\ell_0^\nu\eta^{-1})\eta\sigma_\infty(f^{-1}\ell_0^\nu\eta^{-1})^{-1}.
 \tag{28bis}
\]

Ces formules permettent de mettre en relation, dans le cas $h=1$ les solutions des \deleted{\unsure{deux}} équations
\[
 g=\alpha\rho(\alpha)^{-1}
 \qquad\text{et}\qquad
 \xi=\alpha_0\rho(\alpha_0)^{-1}
 \tag{29}
\]
par
\[
 \alpha=(\ell_0^\nu f)^{-1}\alpha_0=f^{-1}\ell_0^{-\nu}\alpha_0,
 \tag{30}
\]
et dans le cas où $g=1$, les solutions des deux équations
\[
 h=\beta\sigma_\infty(\beta)^{-1},\qquad
 \eta=\beta_0\sigma_\infty(\beta_0)^{-1}
 \tag{31}
\]
par
\[
 \beta=(f^{-1}\ell_0^\nu\eta^{-1})\beta_0.
 \tag{32}
\]

La dissymétrie entre ces deux formules ((32) étant un peu plus compliquée que (30) par le facteur supplémentaire $\eta^{-1}$ du coefficient de $\beta_0$) provient du fait que $f$ a été défini via la \unsure{formule|formules} des lacets (6), première forme, où $\ell_0$ est défini comme le carré de $\sigma_\infty\rho=\varepsilon_0$, et non de $\rho\sigma_\infty=\varepsilon_1$
\margin{143}

\scanpage{145}
\mspage{362}

Si on identifie maintenant $\widehat{\mathcal G}_{0,3}$ à $\mathrm{Gl}(2,\mathbb Z)^{\prime\wedge}$, et qu'on regarde les \deleted{deux} images des éléments $g,h,f$ etc dans le quotient
\[
 \mathrm{Gl}(2,\widehat{\mathbb Z})^{\pm1\prime}
 =\mathrm{Gl}(2,\widehat{\mathbb Z})^{\pm1}/\pm1
\]
--- qu'on désigne par les $\overline{\mathrm m}$s lettres, par abus de langage --- \inserted{on} est amené à calculer avec des éléments de $\mathrm{Gl}(2,\widehat{\mathbb Z})$, plutôt que de $\mathrm{Gl}(2,\widehat{\mathbb Z})'$. On \unsure{choisira} $\alpha,\beta$ dans $\widehat{\Pi}_{0,3}$ et $\widehat{\widetilde{\Pi}}_{0,3}$ conformément à (13), d'où
\[
 \alpha,\beta\in\mathrm{Gl}(2,\widehat{\mathbb Z})^{\pm1}
 \quad\text{\uline{modulo signe}},
 \tag{33}
\]
mais les éléments correspondants
\[
 g=\alpha\rho(\alpha)^{-1},\quad
 h=\beta\sigma_\infty(\beta)^{-1}
 \quad\in\mathrm{Sl}(2,\widehat{\mathbb Z})^{\deleted{\unsure{\pm1}}}
 \tag{34}
\]
sont déterminés dans $\mathrm{Gl}(2,\widehat{\mathbb Z})^{\deleted{\unsure{\pm1}}}$ \uline{sans ambiguïté de signe}. De $\overline{\mathrm m}$
\[
 \overset{\inserted{\text{(cas }h=1\text{)}}}{\alpha_0}
 \;\unsure{\text{et}|,}\;
 \overset{\inserted{\text{(cas }g=1\text{)}}}{\beta_0}
 \in\mathrm{Gl}(2,\widehat{\mathbb Z})^{\pm1}
 \quad\text{modulo signe},
 \tag{35}
\]
et alors
\[
 \begin{aligned}
 \xi&=\alpha_0\rho(\alpha_0)^{-1}
 \inserted{\quad\text{(cas }h=1\text{)}},\quad
 \eta=\beta_0\sigma_\infty(\beta_0)^{-1}
 \inserted{\quad\text{(cas }g=1\text{)}}
 \in\mathrm{Sl}(2,\widehat{\mathbb Z})
 \end{aligned}
 \tag{36}
\]
sont déterminés sans ambiguïté de signe. On a
\deleted{\unsure{Ceci fait|Ceci fixé} \unsure{(34)|(37)}, les éléments}
\deleted{les relations}
\[
 \deleted{\unsure{\varphi=} }f\in\mathrm{Sl}(2,\widehat{\mathbb Z})
 \quad\text{mod \uline{signe}}
 \quad\text{(pour }f\in\widehat{\Pi}_{0,3}\text{ fixé)}
 \tag{37}
\]
et on en déduit
\[
 \varphi=f\sigma_\infty(f)^{-1},\quad
 \psi=f\rho(f)^{-1}\in\mathrm{Sl}(2,\widehat{\mathbb Z})
 \tag{38}
\]
sans ambiguïté de signe.

\deleted{Les relations} Dans le choix des \deleted{\unsure{repr}} représentants

\deleted{\unsure{(39)}\ \illegible{1}\ $\unsure{\varphi|\xi}=$}
\margin{144}

\scanpage{146}

de $\varepsilon_0\unsure{\in}\mathrm{Sl}(2,\widehat{\mathbb Z})'$ dans $\mathrm{Sl}(2,\widehat{\mathbb Z})$ il y a ambiguïté de signe
\[
 \varepsilon_0=\pm\begin{pmatrix}1&-1\\0&1\end{pmatrix},
\]
\margin{encore que le signe $-$ ne serait \unsure{guère} un choix \unsure{raisonnable} --- on aime voir un élément unip. \deleted{\illegible{1}} engendrant un \unsure{ss-groupe} à 1 paramètre\ldots}
mais pour
\[
 \ell_0=\varepsilon_0^2=\begin{pmatrix}1&-2\\0&1\end{pmatrix}
 \tag{39}
\]
il n'y a pas d'ambiguïté de \uline{signe}. Ceci étant, les relations
\[
 \xi=\ell_0^\nu\varphi\qquad\text{(cas }h=1\text{)}
 \tag{40}
\]
et
\[
 \eta=\ell_0^{-\nu}\psi\qquad\text{(cas }g=1\text{)}
 \tag{41}
\]
sont elles vraies sans \uline{signe} correcteur \unsure{??|?!}

Le cas $h=1$ correspond aux données (dans $\widehat{\widetilde{\Pi}}_{0,3}$) de $\xi$ et $\deleted{\unsure{\varphi}}f$, \inserted{et $\nu$,} satisfaisant
\[
 \xi\rho(\xi)\rho^2(\xi)=1,\qquad
 \xi=\ell_0^\nu
 \bigl(\underbrace{f\sigma_\infty(f)^{-1}}_{\varphi}\bigr),
\]
\deleted{et l'équation en $f$ \unsure{est}}
ou encore aux données
\deleted{$\alpha_0\in\widehat{\widetilde{\Pi}}_{0,3},\ f\in\widehat{\Pi}_{0,3},$}
\[
 \alpha_0\in\widehat{\widetilde{\Pi}}_{0,3},\quad
 f\in\widehat{\Pi}_{0,3},\quad\nu\in\widehat{\mathbb Z}
\]
satisfaisant les relations
\[
 \underbrace{\alpha_0\rho(\alpha_0)}_{\xi}
 =\ell_0^\nu\underbrace{f\sigma_\infty(f)^{-1}}_{\varphi}
 \tag{42}
\]
en termes de ces données on définit
\[
 \left\{\begin{aligned}
 \xi&=\alpha_0\rho(\alpha_0)\\
 \varphi&=f\sigma_\infty(f)^{-1}\\
 g&=\sigma_\infty(f)^{-1}\ell_1^\nu\rho(f)
 =(\ell_0^\nu f)^{-1}(\ell_0^\nu\varphi)\rho(\ell_0^\nu f)
 \end{aligned}\right.
 \tag{43}
\]
Ceci rappelé, on remonte $\alpha_0,f$ dans $\mathrm{Gl}(2,\widehat{\mathbb Z})^{\pm1}$,
\margin{145}

\scanpage{147}
\mspage{363}

avec ambiguïté de signe, \deleted{\unsure{mais on}} et \deleted{\illegible{2}}\inserted{$\xi$ et $\varphi$} sont donc déterminés dans $\unsure{\mathrm{Sl}|\mathrm{Gl}}(2,\widehat{\mathbb Z})$, indépendamment des choix de signes, \deleted{\illegible{2}} et de $\overline{\mathrm m}$ $g$.

On \unsure{aura|a} donc dans $\mathrm{Sl}(2,\widehat{\mathbb Z})$
\[
 \deleted{\xi=\alpha_0\rho(\alpha_0)=\xi\,\unsure{\varepsilon\ell_0^\nu}}
\]
\[
 \xi=\varepsilon\ell_0^\nu\varphi
 \qquad\text{i.e.}\qquad
 \alpha_0\rho(\alpha_0)=\varepsilon\ell_0^\nu f\sigma_\infty(f)^{-1},
 \tag{44}
\]
\inserted{\uline{dans $\mathrm{Sl}(2,\widehat{\mathbb Z})$},}
avec un signe $\varepsilon\in\{\pm1\}$ bien déterminé, et la question qui se pose est si on a \unsure{tjrs|tjs} $\varepsilon=+1$.

\deleted{\unsure{Définissant} $f,\unsure{\rho_0|\rho}$}
\inserted{Notons que la formule}
\[
 \rho(\ell_0)=\ell_1
\]
est vraie dans $\mathrm{Sl}(2,\widehat{\mathbb Z})$ (sans \uline{signe} correcteur) donc la \inserted{dernière} expression de $g$ dans (43) est valable dans $\mathrm{Sl}(2,\widehat{\mathbb Z})$ (sans \uline{signe}), ce qui montre que $g\in\mathrm{Sl}(2,\widehat{\mathbb Z})$, \deleted{\unsure{n'est pas}} est un $\rho$-transformé de $\ell_0^\nu\varphi$, \deleted{\unsure{donc}} cette dernière est un $\rho$-transformé de $g$.
\deleted{\illegible{1} $g$, En fait (\illegible{3} $g$) la relation}
\inserted{On a}
\[
 g\rho(g)\rho^2(g)=\varepsilon_0.1
 \qquad\varepsilon\in\{\pm1\},
\]
\[
 \deleted{(\ell_0^\nu\varphi)\rho(\ell_0^\nu\varphi).}
\]
et \deleted{de $\overline{\mathrm m}$} on en déduit, si $g'=\alpha g\rho(\alpha)^{-1}$
\deleted{\unsure{$\rho^2\in$}} ($\alpha\in\mathrm{Gl}(2,\widehat{\mathbb Z})$), qu'on a
\[
 g'\rho(g')\rho^2(g')=\varepsilon_0.1
\]
donc en particulier pour $g'=\ell_0^\nu\varphi\inserted{=\varepsilon\xi\quad\text{(44)}}$ \deleted{d'où}, avec $\varepsilon\in\{\pm1\}$ \deleted{\illegible{2}} or $\xi$ donné par (42) satisfait à
\[
 \xi\rho(\xi)\rho^2(\xi)=1
\]
(\uline{NB} on n'a pas $\rho^3=1$ dans $\mathrm{Sl}(2,\widehat{\mathbb Z})$, mais $\rho^3=-1$, ce qui implique que l'autom. int. \deleted{\unsure{qu'il}} \unsure{défini|définit} par $\rho$ est d'ordre 3\ldots), d'où
\[
 \varepsilon_0=\varepsilon,
\]
\margin{146}

\scanpage{148}

i.e. la question de la validité de la formule
\[
 \xi=\ell_0^\nu\varphi
\]
dans $\mathrm{Sl}(2,\widehat{\mathbb Z})$ équivaut à celle de la validité de la formule
\[
 g\rho(g)\rho^2(g)=1,
\]
que nous laissons en suspens pour le moment. Notons pour mémoire les formules suivantes \uline{dans} \uline{$\unsure{\mathrm{Gl}|\mathrm{Sl}}(2,\widehat{\mathbb Z})$}

\inserted{$\det\alpha_0\in\pm1,\ \det f=1$,}
\[
\begin{gathered}
\text{(Cas }h=1\text{)}\\
\left\{
\begin{aligned}
 &\alpha_0,f\in\mathrm{Gl}(2,\widehat{\mathbb Z})^{\pm1}
 \quad \deleted{\unsure{\text{avec}}}\ \text{(signes arbitraires, on les \uline{choisit})}\\
 &\left.
 \begin{aligned}
 \xi&\overset{\mathrm{dfn}}=\alpha_0\rho(\alpha_0)^{-1}\\
 \varphi&\overset{\mathrm{dfn}}=f\sigma_\infty(f)^{-1}
 \end{aligned}
 \right|\ \in\mathrm{Sl}(2,\widehat{\mathbb Z})\\
 &\bigl[g\overset{\mathrm{dfn}}=\sigma_\infty(f)^{-1}\ell_1^\nu\rho(f)
 =(\ell_0^\nu f)^{-1}(\ell_0^\nu\varphi)\rho(\ell_0^\nu f)\bigr]\\
 &\text{\unsure{d'où|donc}}\quad
 \xi=\deleted{\unsure{\ell_0^{\nu}}}\cdot\varepsilon\ell_0^\nu\varphi
 \qquad\varepsilon\in\{\pm1\}\\
 &\left[\begin{aligned}
 &g\rho(g)\rho^2(g)=\varepsilon.1\\
 &g=\varepsilon.\alpha\rho(\alpha)^{-1}
 \qquad\alpha=(\ell_0^\nu f)^{-1}\alpha_0=f^{-1}\ell_0^{-\nu}\alpha_0
 \end{aligned}\right]
\end{aligned}
\right.
\end{gathered}
\tag{45}
\]

\margin{Il me semble maintenant plus pratique de \uline{définir} $g$ par $g=\alpha\rho(\alpha)^{-1}$, de façon à avoir ainsi $g\rho(g)\rho^2(g)=1$, et de \unsure{placer} le signe $\varepsilon=\varepsilon_2$ dans la formule $g=\varepsilon\sigma_\infty(f^{-1})\ell_1^{\unsure{\nu}}\rho(f)$ --- \unsure{le signe} \illegible{1} \unsure{n'intervient} \illegible{3} \unsure{dans} $\xi=\varepsilon\ell_0^\nu\varphi$.}
\margin{\uline{NB} \unsure{on écr.} \illegible{1}\quad $\varepsilon=\varepsilon_2$}

De façon symétrique
\[
\begin{gathered}
\text{(Cas }g=1\text{)}\\
\left\{
\begin{aligned}
 &\beta_0,f\in\mathrm{Sl}(2,\widehat{\mathbb Z})
 \qquad\text{(modulo signes, on les \uline{choisit})}\\
 &\eta\overset{\mathrm{dfn}}=\beta_0\sigma_\infty(\beta_0)^{-1}\\
 &\psi\overset{\mathrm{dfn}}=\unsure{\psi|f}\rho(f)^{-1}\\
 &h\overset{\mathrm{dfn}}=f^{-1}\ell_0^\nu\sigma_\infty(\rho(f))=
\end{aligned}
\right.
\end{gathered}
\tag{46}
\]

\figure{\begin{tikzpicture}[x=1mm,y=1mm,line width=.35pt,line cap=round]
\draw (0,.1) .. controls (.35,-.08) and (.9,-.06) .. (1.3,.04);
\draw (2.2,.22) .. controls (2.45,.34) and (2.65,.25) .. (2.9,.22);
\draw (4.1,.38) -- (5.4,.55);
\draw (6.7,.57) -- (7.45,.87);
\end{tikzpicture}}
\margin{147}

\scanpage{149}
\mspage{364}

Je vais regarder le cas $h=1$ (cf \unsure{(45)}), en posant
\[
\left\{\begin{aligned}
 \alpha_0&=\begin{pmatrix}a&b\\c&d\end{pmatrix}\\
 f&=\begin{pmatrix}r&s\\t&u\end{pmatrix}
\end{aligned}\right.
\tag{47}
\]
\margin{changement de signes dans $\alpha_0,\unsure{-f|f}$ est \unsure{inessentiel}}
\margin{NB Dans le cas ``arithmétique'' (i.e. $f$ \unsure{provient|provenant} de $\widehat{\Pi}_{0,3}$, $\alpha_0$ de $\widehat{\widetilde{\Pi}}_{0,3}=\widehat{\Pi}_{0,3}.\{1,\tau_\infty\}$, on aura $\alpha_0,f\equiv1\ (2)$ i.e. $a,d,r,u\equiv1\ (2)$, $b,c,s,t\equiv0\ (2)$}
\unsure{si|où} \deleted{\illegible{1}} $a,b,c,d,r,s,t,u\in\widehat{\mathbb Z}$, avec
\[
 \boxed{ab-cd=\delta\in\{\pm1\}}
 \tag{48}
\]
\[
 \boxed{ru-st=1}.
 \tag{49}
\]

Ceci \unsure{posé}
\[
 \xi=\alpha_0\rho(\alpha_0)^{-1}
 =\begin{pmatrix}A&B\\C&D\end{pmatrix}
 \tag{50}
\]
\[
 \varphi=f\underbrace{\sigma_\infty(f)^{-1}}_{{}^{t}f}
 =\begin{pmatrix}R&S\\T&\unsure{U|u}\end{pmatrix}
 =\inserted{\unsure{?}}
 \begin{pmatrix}r^2+s^2&rt+su\\\unsure{rt+su}&t^2+u^2\end{pmatrix}
 \tag{51}
\]
de sorte qu'on aura (cf \S\,42, formule (78)
\[
\left\{\begin{aligned}
 \delta A&=(a^2+b^2+ab)-(ac+bd+bc)
 &\delta B&=ac+bd+bc\\
 \delta C&=ac+bd+ad-(c^2+d^2+cd)
 &\delta D&=c^2+d^2+cd
\end{aligned}\right.,
\tag{52}
\]
\margin{NB on a $\delta=\unsure{\varepsilon_3}$}
et (\S\,42, (90) et 91)
\[
 R=r^2+s^2,\qquad S=T=rt+su,\qquad U=t^2+u^2.
 \tag{53}
\]
\deleted{\illegible{1}}
\deleted{\illegible{2}}\inserted{On aura} les formules
\[
 C+D\unsure{+|-}B=1
 \tag{54}
\]
\[
 S=T.
 \tag{55}
\]
La relation essentielle reliant $\alpha_0$ et $f$, i.e. reliant les ``\unsure{variables}'' $a,b,c,d$ avec $u,r,s,t$, est la relation
\[
 \boxed{\xi=\varepsilon\ell_0^\nu\varphi,\qquad\varepsilon\in\{\pm1\}}
 \tag{56}
\]
\margin{``relation des lacets''}
qui s'écrit aussi $\ell_0^{-\nu}\xi=\varepsilon\varphi$, i.e.
\margin{148}
